\documentclass[10pt,a4paper]{amsart}
\usepackage[margin=19mm]{geometry}
\usepackage{amsmath,amssymb,mathtools}
\usepackage[scr=rsfs,cal=euler]{mathalfa}
\usepackage{xfrac}
\usepackage[unicode,hidelinks,bookmarks=false]{hyperref}
\newcommand{\ie}{i.e.\ }
\newcommand{\cf}{cf.\ }
\newcommand{\resp}{respectively\ }
\newcommand{\Subsection}{\subsection}
\providecommand{\nobreakdash}{\nobreak\hbox{-}}
\setlength{\emergencystretch}{2em}
\multlinegap0pt
\usepackage{enumerate}
\usepackage{amssymb}
\usepackage{mathtools}
\usepackage[shortlabels]{enumitem}
\usepackage{float}
\usepackage{stackengine}
\usepackage{xcolor}
\newtheorem{theorem}{Theorem}
\newcommand{\TT}{\mathbb T}
\title{Asymptotic models for viscoelastic one-dimensional blood flow}
\author{Diego Alonso-Or\'an, Rafael Granero-Belinch\'on, Carlos Yanes P\'erez}
\date{Source: arXiv:2604.05679v1 (CC BY 4.0)}
\begin{document}
\maketitle

\noindent\emph{Source and licence.} Asymptotic models for viscoelastic one-dimensional blood flow. Version arXiv:2604.05679v1 (CC BY 4.0), distributed by arXiv under the Creative Commons Attribution 4.0 International licence, http://creativecommons.org/licenses/by/4.0/, as recorded in the arXiv licence metadata for this exact version and shown on the abstract page https://arxiv.org/abs/2604.05679v1. Full licence notice: \texttt{licenses/arxiv-2604.05679-CC-BY-4.0.txt}.\par\smallskip

\noindent\textit{Editorial scope.} Counted unit: the article's theorem th:bbm-global (source0.tex lines 1522--1534). The article's complete original proof of it is delivered with it (source0.tex lines 1536--1680, one \texttt{proof} environment of 145 lines). No mathematics is written, repaired or re-derived: the statement and the proof are the article's own lines, in the article's own notation, and no retained line is substituted or re-flowed.  Retained uncounted as the source-local context the proof needs: lines 275--283; lines 850--856.  Nothing was omitted from any retained span, so the package carries no omission record.  No retained line contains a citation, so the wrapper carries no bibliography and no bibliography entry.  No retained line contains a figure environment, an image-inclusion command or drawing code, so no image is delivered and the package declares no asset dependency.  Editorial formatting only, in addition: an AMS \texttt{amsart} preamble that restates the article's own declarations and macros used by the retained lines, namely the environments \texttt{theorem} on separate counters, together with the standard packages those lines need.  The retained environments print in this document as Theorem 1 in order of appearance, whereas the article prints the same items with the numbering of its own full document.  The complete native source of the article as distributed in the arXiv e-print for this exact version is bundled unchanged as \texttt{sources/197974/source0.tex}.
\begin{align}\label{asymptotic:model:contr}
	f_{t}&=\mathcal{M}\mathcal{P}\bigg[ - \frac{1}{\varepsilon}\Big(1-\frac{\beta}{2}\Big)f_{xx}
	+ \frac{1}{\varepsilon}\kappa f_{x}
	- \frac{1}{\varepsilon}\frac{\nu}{2}f_{xxx}
	+ \Big(2+\frac{\beta}{4}\Big)(f f_{x})_{x}    \nonumber \\
	&\hspace{4.7cm}+ \frac{1}{\varepsilon}\frac{\nu}{4}f_{x}f_{xx}
	-\frac{\nu}{4}ff_{xxx}
	-2\kappa ff_{x}\bigg],
\end{align}

\begin{theorem}\label{th:existence:strong}
	Let $\nu>0$ and $\kappa>0$, let $s>\frac{5}{2}$, and let $f_0\in H^s(\TT)$ be mean-zero initial data for \eqref{asymptotic:model:contr}.
	Then there exists $T>0$, depending only on $\|f_0\|_{H^s(\TT)}$ and the parameters of the model, and a unique strong solution
	\[
	f\in C([0,T];H^s(\TT)),\qquad f(\cdot,0)=f_0.
	\]
\end{theorem}

\begin{theorem}\label{th:bbm-global}
	Assume $\nu=0$ and $\beta>-2$. Let $f_0\in H^2(\TT)$ be mean-zero and let
	$f\in C([0,T_{\max});H^2(\TT))$ be the maximal strong solution of \eqref{asymptotic:model:contr} with $f(0)=f_0$. Then there exists $\delta_0=\delta_0(\varepsilon,\kappa,\beta)>0$ such that if
	\[
	\|f_0\|_{H^2}\le \delta_0,
	\]
	the solution is global, $T_{\max}=+\infty$, and there exist constants $C,c>0$, depending only on
	$(\varepsilon,\kappa,\beta)$, such that
	\[
	\|f(t)\|_{H^2}\le C e^{-ct}\|f_0\|_{H^2}\qquad\text{for all }t\ge 0.
	\]
	In particular, $\|f(t)\|_{H^2}\to 0$ as $t\to\infty$.
\end{theorem}

\begin{proof}[Proof of Theorem \ref{th:bbm-global}]
	We only establish the a prior energy estimates leading to global control and decay under the smallness
	assumption on $\|f_0\|_{H^2}$. The construction of solutions via the approximation/compactness procedure follow by the same standard steps used in the local well-posedness theorem
	(Theorem~\ref{th:existence:strong}). To avoid repetition we omit these routine details. 
	\medskip
	
	First notice that when $\nu=0$, the operators $\mathcal{P}$ and $\mathcal{M}$ reduce to
	\[
	\mathcal{P}=\frac{1}{\kappa}\,\mathrm{Id},\qquad
	\mathcal{M}=\Big(\mathrm{Id}-\frac{4}{\kappa^2}\partial_{xx}\Big)^{-1}\Big(\mathrm{Id}+\frac{2}{\kappa}\partial_x\Big),
	\]
	and therefore \eqref{asymptotic:model:contr} can be written in the local BBM-type form
	\begin{equation}\label{eq:bbm:general}
		f_t-\frac{4}{\kappa^2}\partial_{xx}f_t
		=\Big(1+\frac{2}{\kappa}\partial_x\Big)\Big[
		-\frac{1}{\varepsilon}\Big(1-\frac{\beta}{2}\Big)f_{xx}
		+\frac{\kappa}{\varepsilon}f_x
		+\Big(2+\frac{\beta}{4}\Big)(ff_x)_x
		-2\kappa ff_x\Big].
	\end{equation}
	We set
	\[
	a:=\frac{1}{\varepsilon}\Big(\frac{\beta}{2}-1\Big),\qquad
	b:=\frac{\kappa}{\varepsilon},\qquad
	c:=2+\frac{\beta}{4},\qquad
	d:=2\kappa,
	\]
	so that the bracket in \eqref{eq:bbm:general} is $Q:=a f_{xx}+b f_x+c(ff_x)_x-dff_x$. \medskip
	
	We introduce the following energies
	\[
	E_1(t):=\|f(t)\|_{L^2}^2+\frac{4}{\kappa^2}\|f_x(t)\|_{L^2}^2,
	\qquad
	E_2(t):=\|f_x(t)\|_{L^2}^2+\frac{4}{\kappa^2}\|f_{xx}(t)\|_{L^2}^2,
	\]
	and the dissipations $D_1(t):=\|f_x(t)\|_{L^2}^2$, $D_2(t):=\|f_{xx}(t)\|_{L^2}^2$. Note that $E_1\simeq\|f\|_{H^1}^2$ and
	$E_2\simeq\|f\|_{H^2}^2$ (with constants depending only on $\kappa$). Since $f$ has zero mean, Poincar\'e inequality yields
	\begin{equation}\label{eq:coercive-bbm}
		E_1(t)+E_2(t)\ \lesssim\ D_1(t)+D_2(t).
	\end{equation}
	We will also use the one-dimensional Gagliardo--Nirenberg bounds
	\begin{equation}\label{eq:GN-BBM}
		\|f\|_{L^\infty}\lesssim \|f\|_{L^2}^{1/2}\|f_x\|_{L^2}^{1/2}\lesssim \sqrt{E_1},
		\qquad
		\|f_x\|_{L^\infty}\lesssim \|f_x\|_{L^2}^{1/2}\|f_{xx}\|_{L^2}^{1/2}\lesssim \sqrt{E_2}.
	\end{equation}
	
	Taking the $L^2$ inner product of \eqref{eq:bbm:general} with $f$ and integrating by parts gives
	\[
	\frac12\frac{d}{dt}\|f\|_{L^2}^2-\frac{4}{\kappa^2}\int_\TT f_{xxt}f\,dx
	=\frac12\frac{d}{dt}\|f\|_{L^2}^2+\frac{2}{\kappa^2}\frac{d}{dt}\|f_x\|_{L^2}^2
	=\frac12\frac{d}{dt}E_1(t),
	\]
	while on the right-hand side
	\[
	\int_\TT\Big(1+\frac{2}{\kappa}\partial_x\Big)Q\,f\,dx
	=\int_\TT Q f\,dx-\frac{2}{\kappa}\int_\TT Q f_x\,dx.
	\]
	The linear contributions are
	\[
	\int_\TT a f_{xx}f\,dx=-aD_1,\qquad
	-\frac{2}{\kappa}\int_\TT b f_x f_x\,dx=-\frac{2b}{\kappa}D_1,
	\]
	whereas $\int_\TT b f_x f\,dx=0$ and $-\frac{2}{\kappa}\int_\TT a f_{xx} f_x\,dx=0$. Thus the linear part yields
	$-(a+2b/\kappa)D_1$. For the nonlinear part, using
	\[
	\int_\TT (ff_x)_x f\,dx=-\int_\TT f\,f_x^2\,dx,\qquad
	\int_\TT (ff_x)_x f_x\,dx=\frac12\int_\TT f_x^3\,dx,\qquad
	\int_\TT f f_x f_x\,dx=\int_\TT f\,f_x^2\,dx,
	\]
	we obtain
	\[
	\Big|\int_\TT (c(ff_x)_x-dff_x)f\,dx\Big|\lesssim \|f\|_{L^\infty}D_1,\qquad
	\Big|\int_\TT (c(ff_x)_x-dff_x)f_x\,dx\Big|\lesssim \|f_x\|_{L^\infty}D_1.
	\]
	Invoking \eqref{eq:GN-BBM} we conclude that
	\begin{equation}\label{ineq:E1}
		\frac12\frac{d}{dt}E_1(t)+\Big(a+\frac{2b}{\kappa}\Big)D_1(t)
		\le C\big(\sqrt{E_1(t)}+\sqrt{E_2(t)}\big)D_1(t),
	\end{equation}
	for some $C=C(\varepsilon,\kappa,\beta)$. \medskip
		
	Next, taking the $L^2$ inner product of \eqref{eq:bbm:general} with $-f_{xx}$ and integrating by parts yields
	\[
	\frac12\frac{d}{dt}\|f_x\|_{L^2}^2-\frac{4}{\kappa^2}\int_\TT f_{xxt}(-f_{xx})\,dx
	=\frac12\frac{d}{dt}\|f_x\|_{L^2}^2+\frac{2}{\kappa^2}\frac{d}{dt}\|f_{xx}\|_{L^2}^2
	=\frac12\frac{d}{dt}E_2(t),
	\]
	and
	\[
	\int_\TT\Big(1+\frac{2}{\kappa}\partial_x\Big)Q\,(-f_{xx})\,dx
	=\int_\TT Q(-f_{xx})\,dx-\frac{2}{\kappa}\int_\TT Q(-f_{xxx})\,dx.
	\]
	The linear terms satisfy
	\[
	\int_\TT a f_{xx}(-f_{xx})\,dx=-aD_2,\qquad
	-\frac{2}{\kappa}\int_\TT b f_x(-f_{xxx})\,dx
	=\frac{2b}{\kappa}\int_\TT f_x f_{xxx}\,dx=-\frac{2b}{\kappa}D_2,
	\]
	while $\int_\TT b f_x(-f_{xx})\,dx=0$ and $-\frac{2}{\kappa}\int_\TT a f_{xx}(-f_{xxx})\,dx=0$, hence the linear part yields
	$-(a+2b/\kappa)D_2$. For the nonlinear terms we use we find that
	\[
	\Big|\int_\TT (ff_x)_x f_{xx}\,dx\Big|
	\lesssim \|f_x\|_{L^\infty}\|f_x\|_{L^2}\|f_{xx}\|_{L^2}+\|f\|_{L^\infty}\|f_{xx}\|_{L^2}^2
	\lesssim (\|f\|_{L^\infty}+\|f_x\|_{L^\infty})D_2,
	\]
	and similarly
	\[
	\Big|\int_\TT (ff_x)_x f_{xxx}\,dx\Big|
	+\Big|\int_\TT (f f_x) f_{xxx}\,dx\Big|
	\lesssim (\|f\|_{L^\infty}+\|f_x\|_{L^\infty})D_2.
	\]
	Using \eqref{eq:GN-BBM} and Young's inequality to absorb mixed products we obtain
	\begin{equation}\label{ineq:E2}
		\frac12\frac{d}{dt}E_2(t)+\Big(a+\frac{2b}{\kappa}\Big)D_2(t)
		\le C\big(\sqrt{E_1(t)}+\sqrt{E_2(t)}\big)D_2(t).
	\end{equation}
	
	Setting $F:=E_1+E_2$, $G:=D_1+D_2$ and $\mu:=a+\frac{2b}{\kappa}>0$, adding \eqref{ineq:E1} and \eqref{ineq:E2} gives
	\begin{equation}\label{eq:F-ineq}
		\frac12 F'(t)+\mu\,G(t)\le C\sqrt{F(t)}\,G(t).
	\end{equation}
	Choose $\delta_0>0$ such that $C\delta_0\le \mu/2$ and assume $\|f_0\|_{H^2}\le\delta_0$ (absorbing the equivalence between
	$\|f_0\|_{H^2}$ and $F(0)^{1/2}$ into $\delta_0$). Define
	\[
	\mathcal{I}:=\Big\{t\in[0,T_{\max})\,:\,F(s)\le \delta_0^2\ \text{for all }s\in[0,t]\Big\},\qquad
	T_*:=\sup\mathcal{I}.
	\]
	Then $T_*>0$ and $F(t)\le \delta_0^2$ for $t\in[0,T_*)$. On this interval \eqref{eq:F-ineq} implies
	\[
	\frac12 F'(t)+\mu G(t)\le C\delta_0 G(t)\le \frac{\mu}{2}G(t),
	\qquad\text{hence}\qquad
	F'(t)+\mu G(t)\le 0 \quad \text{on }[0,T_*).
	\]
	In particular $F$ is non-increasing on $[0,T_*)$ and thus $F(t)\le F(0)\le \delta_0^2$ there; by continuity this forces
	$T_*=T_{\max}$, so the inequality $F'(t)+\mu G(t)\le 0$ holds for all $t\in[0,T_{\max})$.
	
	Finally, using the coercivity \eqref{eq:coercive-bbm} we have $G\gtrsim F$, hence
	\[
	F'(t)+c_1 F(t)\le 0
	\]
	for some $c_1=c_1(\varepsilon,\kappa,\beta)>0$, which yields $F(t)\le F(0)e^{-c_1 t}$. Since $F\simeq \|f\|_{H^2}^2$,
	we conclude $\|f(t)\|_{H^2}\le C e^{-ct}\|f_0\|_{H^2}$ for suitable $C,c>0$, and in particular $T_{\max}=+\infty$ and
	$\|f(t)\|_{H^2}\to 0$ as $t\to\infty$.
\end{proof}

\end{document}
